\documentclass[11pt]{article}
\usepackage{booktabs}
\usepackage[hidelinks]{hyperref}
\title{Continuing after a motion guard rejects a move}
\date{}
\begin{document}
\maketitle

A motion guard can prevent a collision without telling the scheduler
how to finish outstanding work. In the previous MRWS audit, the guard
stopped the simulation on a request to enter an occupied cell. We
compared that behavior with two ways of continuing at the same guard.
Retrying the old path removed the terminal abort but still stalled
most runs. Invalidating that path allowed more orders to finish,
without improving every run or guaranteeing recovery.

\section{Matched development campaign}
Use the same public MRWS source, layout, non-genetic native
\texttt{simple-interrupt} scheduler and operating scope as the
\href{https://github.com/eternal-november/eos-research-programmes/pull/401}
{source audit}: the \(7\times11\) four-robot layout, capacity three,
and horizons of 240 ticks. The development protocol reused the
16 development seeds; it did not select an unseen case.
Actuator and sensor faults were enabled at \(1/64\) and \(1/128\)
per applicable draw; battery faults were off. Offered orders have
a separate deterministic generator. A fixed table supplies four
opportunity draws per tick and actor; the native fault routine
still decides whether a draw can realize a new fault. Actual
faults can differ when histories change. Merely reusing a global
seed would not match the exogenous inputs because branches use
random numbers differently. The replay checks the same opportunity
table in each group, the same four offered orders with their
contents, priority and release times in the two continuing arms,
and the corresponding prefix in the terminal arm.

At the native move operation, keep the occupied-cell check and
its ordering before the position update. Compare three dispositions
when it fails:
\begin{description}
\item[Native stop.] Run the original guard, which ends the run before
  applying the rejected move.
\item[Wait.] End only the current action and keep the cached path
  ready for an ordinary later retry.
\item[Replan.] End the current action and discard the blocked actor's
  cached path. On its next applicable turn normal native planning
  can try a new path.
\end{description}
Do not change location, carried inventory, order authority or
service credit on refusal. The harness checks the first two
unchanged local state fields around each refusal and calls the
original operation for allowed moves. A refusal inside a nested
native move still applies to that actual moving robot's cached
path. Earlier lawful effects in the tick are not undone.
Both continuing arms act on the same local physical rejection and
use the ordinary native planner on its existing information.
No globally hidden fault variable or future motion is passed to
either as a new input. These are comparisons inside the simulator's
own routing information, not claims about data available to a
decentralized robot in a real warehouse.

\section{Results}
\begin{center}
\begin{tabular}{lrrrl}
\toprule
Policy & Orders completed & Runs completed & Rejections & Other endings\\
\midrule
Native stop & 2 & 0 & 16 & 16 guard stops\\
Wait & 34 & 1 & 1610 & 15 at time limit\\
Replan & 49 & 5 & 62 & 11 at time limit\\
\bottomrule
\end{tabular}
\end{center}
The two continuing arms had no native occupied-cell aborts,
and all attained tick-boundary configurations checked had
distinct positions. Replanning finished more orders than
waiting on nine paired runs, the same on five, and fewer
on two. This is a finite comparison for the pinned plant and
control logic, not a distributional claim or proof of dominance.
Orders count actual native station completion, not acknowledged
intent. Rejections are events, not distinct safety incidents.
Waiting can retry the same blocked move on many ticks.
A stopped run is not coded as a useful safe completion and no
timed-out run was dropped from the comparison.

\section{What the comparison settles}
Refusal should leave the plant and its existing work intact
rather than make the rejected effect appear applied. That
change is necessary to compare progress past the original
guard. Dropping the cached path then materially helped in this
development campaign, but even ordinary replanning did not
recover all orders and sometimes did worse than waiting.
Neither public knowledge that a move was rejected nor retention
of carried load establishes a route or a deadline guarantee.
A later transfer must expose the exact earlier duty that remains
live and the independently observed event that settles it.

This is the direct guarded-replanning baseline for this source,
not an improvement over it, an independent transfer, or an
architecture acceptance claim. The model tests counted inventory,
not unique-item custody or field safety.
\end{document}

